% Kintsugi -- Relazione di sintesi
% Versione breve del documento di architettura (main.tex), pensata per la lettura dei docenti.
% Compilare con: make sintesi   (oppure: latexmk -pdf main-sintesi.tex)

\documentclass[11pt, a4paper, oneside, openany]{report}

% Pacchetti del documento. Lo stile grafico è in config/stile.tex,
% i metadati (titolo, autori, versione) in config/metadati.tex.

\usepackage[utf8]{inputenc}
\usepackage[T1]{fontenc}
\usepackage[italian]{babel}
\usepackage{lmodern}
\usepackage{microtype}

\usepackage[a4paper, top=2.8cm, bottom=2.8cm, inner=3cm, outer=2.5cm, headheight=14pt]{geometry}
\usepackage{fancyhdr}
\usepackage{titlesec}

\usepackage{xcolor}
\usepackage{graphicx}
\usepackage{tikz}
\usetikzlibrary{positioning, arrows.meta, fit, backgrounds, calc, shapes.geometric, shapes.multipart, shadows, decorations.pathreplacing}

\usepackage{booktabs}
\usepackage{array}
\usepackage{tabularx}
\usepackage{longtable}
\usepackage{enumitem}
\usepackage{caption}
\usepackage{amsmath}
\usepackage{tcolorbox}

\usepackage{hyperref}

% Valori numerici generati da config/parametri.yaml (li crea make): si usano con \parametro{chiave}.
\IfFileExists{config/parametri.tex}{\input{config/parametri}}{%
  \PackageError{parametri}{Manca config/parametri.tex}{Lancia make nella cartella docs: crea il file da config/parametri.yaml.}}

\usepackage{float}
% Metadati del documento: si modificano solo qui.

\newcommand{\progetto}{Kintsugi}
\newcommand{\titolodoc}{Documento di architettura}
\newcommand{\sottotitolo}{Laboratorio di affidabilità e asincronia su Kubernetes e AWS}
\newcommand{\versionedoc}{0.10}
\newcommand{\statodoc}{Bozza}

\newcommand{\universita}{Università degli Studi di Napoli Federico II}
\newcommand{\scuola}{Scuola Politecnica e delle Scienze di Base}
\newcommand{\dipartimento}{Dipartimento di Ingegneria Elettrica e delle Tecnologie dell'Informazione}
\newcommand{\corso}{Corso di Laurea Magistrale in Ingegneria Informatica}
\newcommand{\insegnamento}{Software Project Management \& Evolution}
\newcommand{\docenti}{Prof.\ Sergio Di Martino \\ Prof.\ Luigi Libero Lucio Starace}
\newcommand{\annoaccademico}{2025/2026}

\newcommand{\autoreuno}{Francesco Terrecuso}
\newcommand{\matricolauno}{DE5000005}
\newcommand{\autoredue}{Gianmarco Vellotti}
\newcommand{\matricoladue}{DE5000041}


% Metadati propri della sintesi (il documento tecnico mantiene i suoi).
\renewcommand{\titolodoc}{Relazione di sintesi}
\renewcommand{\versionedoc}{0.1}
\renewcommand{\statodoc}{Bozza}

\input{config/stile}

% La sintesi usa le sezioni come livello principale: niente capitoli, niente pagine vuote.
\renewcommand{\thesection}{\arabic{section}}
\renewcommand{\sectionmark}[1]{\markboth{#1}{}}
\titlespacing*{\section}{0pt}{16pt}{6pt}
\makeatletter
\renewcommand\tableofcontents{\section*{Indice}\@starttoc{toc}}
\makeatother

\begin{document}

\pagenumbering{roman}
\input{frontmatter/frontespizio}
\pagenumbering{arabic}

\tableofcontents

\section{Il progetto in breve}
\label{s:progetto}

\progetto{} è un servizio di ordini volutamente minimo: riceve un ordine e lo salva in un database. Lo realizziamo in \textbf{cinque architetture} (V1--V5), gli facciamo accadere \textbf{guasti controllati} (F1--F7) mentre è sotto carico e misuriamo come si ripristina. Non valutiamo il codice, ma le \textbf{scelte architetturali}: ridondanza, coda asincrona, replica dei dati, elasticità.

Il nome viene dal \textit{kintsugi}, l'arte giapponese di riparare la ceramica rotta con l'oro, mettendo in evidenza le crepe: ogni guasto misurato mostra dove l'architettura cede e come si ricompone.

\paragraph{Perché è un tema di architettura.} Nel corso (T06) l'architettura è l'insieme delle decisioni principali, quelle difficili da cambiare in seguito, e determina proprietà come disponibilità, scalabilità e prestazioni. Qui prendiamo quelle decisioni, le mettiamo alla prova con i guasti e ne misuriamo l'effetto, invece di limitarci a descriverlo.

\subsection{Le domande di ricerca}

\begin{table}[H]
\centering
\caption{Le sei domande di ricerca e le versioni confrontate.}
\label{tab:rq}
\begin{tabularx}{\linewidth}{@{}lLl@{}}
\toprule
\textbf{ID} & \textbf{Domanda} & \textbf{Confronto} \\
\midrule
RQ1 & Quanto fermo si evita passando da una a più copie dell'API? & V1 / V2 \\
RQ2 & Cosa cambia, in tempi, ordini persi e ripristino, inserendo una coda tra API e database? & V2 / V3 \\
RQ3 & Quanto dipende il ripristino dal tipo di guasto e dai parametri delle sonde? & tutte \\
RQ4 & Quanto fermo evita una replica del database con failover, e quanti dati si rischiano? & V2 / V4 \\
RQ5 & L'autoscaling sulla lunghezza della coda assorbe meglio un picco, a parità di costo? & V3 / V5 \\
RQ6 & Quali risultati locali si ripetono su AWS, e perché gli altri cambiano? & locale / AWS \\
\bottomrule
\end{tabularx}
\end{table}

\subsection{Il metodo}

Tre regole guidano tutto il lavoro.

\begin{enumerate}
  \item \textbf{Un fattore alla volta.} Ogni architettura cambia un solo aspetto rispetto a quella da cui deriva, così ogni differenza misurata ha una causa attribuibile.
  \item \textbf{Prima prevedere, poi misurare.} Le previsioni si scrivono e si registrano su Git \emph{prima} degli esperimenti (tag \texttt{m1}). Lo scarto tra previsto e misurato mostra quanto abbiamo capito del sistema.
  \item \textbf{Dichiarare i limiti.} Ciò che gli esperimenti non possono dimostrare è scritto esplicitamente (sezione~\ref{s:limiti}).
\end{enumerate}

\begin{figure}[H]
\centering
\resizebox{\linewidth}{!}{%
\begin{tikzpicture}[node distance=0.55cm,
  fase/.style={box, minimum width=2.45cm, minimum height=1.25cm},
  doc/.style={etichetta, text width=2.45cm}]
\node[fase] (p) {\textbf{Prevedere}\\{\scriptsize modelli e parametri}};
\node[fase, right=of p] (c) {\textbf{Configurare}\\{\scriptsize una versione V1--V5}};
\node[fase, right=of c, guastato] (g) {\textbf{Guastare}\\{\scriptsize un componente}\\{\scriptsize (guasti F1--F7)}};
\node[fase, right=of g] (m) {\textbf{Misurare}\\{\scriptsize client, DB, cluster}};
\node[fase, right=of m] (k) {\textbf{Confrontare}\\{\scriptsize previsto e misurato}};
\foreach \a/\b in {p/c, c/g, g/m, m/k} \draw[freccia] (\a) -- (\b);
\node[doc, below=2mm of p] {sezione~\ref{s:previsioni}\\tag \texttt{m1}};
\node[doc, below=2mm of c] {sezione~\ref{s:architetture}};
\node[doc, below=2mm of g] {sezione~\ref{s:guasti}};
\node[doc, below=2mm of m] {sezione~\ref{s:misura}};
\node[doc, below=2mm of k] {report finale\\(RQ1--RQ6)};
\draw[risposta] (k.north) -- ++(0,0.7) -| node[etichetta, pos=0.25, above] {lo scarto spiega cosa non avevamo capito del sistema} (p.north);
\end{tikzpicture}}
\caption{Il metodo: ogni esperimento confronta una previsione scritta in anticipo con una misura.}
\label{fig:metodo}
\end{figure}

\subsection{I concetti del corso che mettiamo alla prova}

\begin{table}[H]
\centering
\caption{Dalla teoria dell'insegnamento agli esperimenti.}
\label{tab:corso}
\begin{tabularx}{\linewidth}{@{}>{\raggedright\arraybackslash}p{4.6cm}L@{}}
\toprule
\textbf{Concetto (lezioni)} & \textbf{Dove lo misuriamo} \\
\midrule
Scalabilità orizzontale e bilanciamento del carico (T06) & V1 contro V2: tre copie dell'API dietro un bilanciatore (RQ1). Un bilanciatore unico è un punto unico di guasto: nel nostro caso il control plane, che non viene mai guastato (sezione~\ref{s:limiti}). \\
Disponibilità e ``numero di nove'' (T06) & Obiettivo del 99,9\% e budget di errore: ogni fermo misurato si traduce in guasti sopportabili al mese (sezione~\ref{s:previsioni}). \\
Failover attivo-attivo e attivo-passivo; \textit{heartbeat} e osservabilità (T06) & Attivo-attivo: le copie dell'API in V2. Attivo-passivo: primario e standby del database in V4. I controlli periodici sono le sonde di Kubernetes, che misuriamo con il blocco di un'istanza (RQ3). \\
Replica del database e suo costo di sincronizzazione (T06) & V4a con replica asincrona (possibile perdita di dati) contro V4s con replica sincrona (scritture più lente) (RQ4). \\
Interazione sincrona e asincrona, coda e broker (T06) & V2 contro V3: la coda separa chi produce gli ordini da chi li salva, e si riempie se i produttori superano i consumatori (RQ2). \\
Scalabilità contro elasticità (T12) & V5: i worker seguono la domanda con una regola dinamica sulla lunghezza della coda; il costo si misura in risorse usate nel tempo, come nel pagamento a consumo del cloud (RQ5). \\
Cloud IaaS: macchine virtuali e bilanciatore gestiti (T12) & Ripetizione di una parte degli esperimenti su AWS (EKS, EC2, ALB) per capire cosa si trasferisce (RQ6). \\
\bottomrule
\end{tabularx}
\end{table}

\paragraph{Come leggere questo documento.} È la versione breve del \emph{Documento di architettura} (\texttt{docs/main.pdf}), che resta il riferimento tecnico completo: contiene le diciassette schede di decisione, i parametri, tutti i diagrammi e i calcoli. I codici D\textit{nnnn} rimandano a quelle schede.

\section{Le cinque architetture}
\label{s:architetture}

Le cinque versioni sono configurazioni dello stesso codice e della stessa immagine (D0012): cambiano solo impostazioni e manifest. Ognuna deriva da un'altra modificando \textbf{un solo fattore} (figura~\ref{fig:derivazione}).

\begin{figure}[H]
\centering
\begin{tikzpicture}
\node[box] (v1) {V1\\{\scriptsize 1 istanza, sincrona}};
\node[box, right=1.3cm of v1] (v2) {V2\\{\scriptsize 3 repliche}};
\node[box, above right=0.2cm and 1.6cm of v2] (v3) {V3\\{\scriptsize asincrona}};
\node[box, below right=0.2cm and 1.6cm of v2] (v4) {V4\\{\scriptsize DB replicato}};
\node[box, right=1.3cm of v3] (v5) {V5\\{\scriptsize autoscaling}};
\draw[freccia] (v1) -- node[above, font=\scriptsize] {+ repliche} (v2);
\draw[freccia] (v2) -- node[above, sloped, font=\scriptsize] {+ coda} (v3);
\draw[freccia] (v2) -- node[below, sloped, font=\scriptsize] {+ replica del DB} (v4);
\draw[freccia] (v3) -- node[above, font=\scriptsize] {+ KEDA} (v5);
\end{tikzpicture}
\caption{Derivazione delle versioni: ogni freccia cambia un solo fattore.}
\label{fig:derivazione}
\end{figure}

La figura~\ref{fig:panoramica} riunisce i componenti. Nelle versioni sincrone (V1, V2, V4) l'API scrive nel database e risponde solo dopo il salvataggio; in quelle asincrone (V3, V5) pubblica l'ordine su una coda (RabbitMQ) e risponde appena la coda lo ha conservato, mentre i worker lo salvano dopo.

\begin{figure}[H]
\centering
\resizebox{\linewidth}{!}{%
\begin{tikzpicture}[every node/.append style={font=\small}]
\node[box, minimum width=1.7cm] (cl) at (0,0) {Cliente\\{\scriptsize JMeter}};
\node[kube] (svc) at (2.6,0) {Service};
\node[copie] (api) at (5.2,0) {API\\{\scriptsize $\times N$}};
\node[db] (db) at (12.0,0.4) {PostgreSQL};
\node[db] (st) at (15.4,0.4) {Standby\\{\scriptsize solo V4}};
\node[coda, label={[etichetta]above:RabbitMQ}] (q) at (9.0,-2.3) {};
\node[copie] (wk) at (12.0,-2.3) {Worker\\{\scriptsize $\times M$}};
\node[kube] (keda) at (9.0,-4.2) {KEDA\\{\scriptsize solo V5}};
\draw[freccia] (cl) -- (svc);
\draw[freccia] (svc) -- (api);
\draw[freccia] (api.east) -- node[etichetta, above] {V1, V2, V4: scrive nel database\\e risponde 201} (db.west);
\draw[freccia] (api.south) |- node[etichetta, pos=0.97, above left] {V3, V5: pubblica\\e risponde 202} (q.west);
\draw[freccia] (q) -- (wk);
\draw[freccia] (wk.north) -- (db.south);
\draw[freccia, double, -{Latex}] (db) -- node[etichetta, above] {replica} (st);
\draw[sonda] (q.south) -- node[etichetta, right] {lunghezza\\della coda} (keda.north);
\draw[sonda] (keda.east) -| node[etichetta, pos=0.25, above] {da 1 a 10 worker} (wk.south);
\end{tikzpicture}}
\caption{I componenti e i due percorsi di un ordine. Le linee tratteggiate sono le richieste di controllo di KEDA e le sue azioni sui worker.}
\label{fig:panoramica}
\end{figure}

\begin{table}[H]
\centering
\caption{Le cinque versioni e i loro punti unici di guasto.}
\label{tab:versioni}
\begin{tabularx}{\linewidth}{@{}lLL@{}}
\toprule
\textbf{Versione} & \textbf{Come funziona} & \textbf{Cosa ferma il servizio} \\
\midrule
V1 & Una sola istanza dell'API. Salva e risponde 201; se non riesce entro 1,5~s risponde 503. & API, database \\
V2 & Tre repliche dell'API. Il Service inoltra solo alle repliche pronte, quindi una replica guasta viene esclusa. & Database \\
V3 & L'API pubblica sulla coda e risponde 202 dopo la conferma del broker; tre worker salvano nel database e confermano solo dopo il salvataggio. & Broker (accettazione); broker e database (completamento: gli ordini attendono in coda) \\
V4 & Come V2, con PostgreSQL primario e standby gestiti da CloudNativePG. V4a replica in modo asincrono, V4s in modo sincrono. & V4a: nessuno, ma gli ultimi ordini possono andare persi. V4s: lo standby, se cade le scritture possono bloccarsi\textsuperscript{*} \\
V5 & Come V3, con KEDA che porta i worker da 1 a 10 in base alla lunghezza della coda. Si confronta anche con V3 a 1, a 3 e a 10 worker fissi. & Come V3 \\
\bottomrule
\end{tabularx}

\vspace{2pt}
{\footnotesize\textsuperscript{*} Comportamento da verificare sulla versione di CloudNativePG adottata (D0014).\par}
\end{table}

\paragraph{Il tempo di un ordine.} In ogni versione si registrano due istanti, entrambi presi nel cluster per evitare orologi diversi (D0005): la ricezione da parte dell'API e la scrittura nel database. La loro differenza è il \emph{tempo end-to-end}, la metrica comune a tutte le versioni. Con la coda il client attende meno, ma l'ordine viene salvato più tardi.

\subsection{Dove si esegue}

Una campagna completa si svolge in \textbf{locale}, su \textbf{kind}, con tre nodi simulati da container Docker (un control plane che non viene mai guastato e due worker, uno per i dati e uno per il calcolo), Kubernetes 1.37 e misure ufficiali su un solo computer (D0013). Una parte degli esperimenti (V2 e V3, carichi W1 e W2, guasti F1, F2 e F3) si ripete su \textbf{AWS}, con EKS, nodi EC2 in tre zone, bilanciatore ALB e dischi EBS, nella regione di Irlanda (D0009). Il costo previsto è di circa 15--20 dollari, con allarme di budget attivo prima di ogni sessione e distruzione del cluster al termine.

Il limite principale dell'ambiente locale è che i nodi di kind sono container sullo stesso computer: la perdita di un nodo è simulata, e per questo la ripetizione su AWS serve a RQ6.

\section{I guasti}
\label{s:guasti}

Ogni run inietta \textbf{un solo guasto}, sotto carico e in un istante registrato, e osserva il ripristino automatico del sistema: nessun intervento manuale.

\begin{table}[H]
\centering
\caption{I sette guasti.}
\label{tab:guasti}
\begin{tabularx}{\linewidth}{@{}lLlll@{}}
\toprule
\textbf{ID} & \textbf{Guasto} & \textbf{Meccanismo} & \textbf{Versioni} & \textbf{RQ} \\
\midrule
F1 & Arresto improvviso di un'istanza dell'API & \texttt{SIGKILL} & V1, V2, V3 & 1, 2, 3 \\
F2 & Arresto di tutti i worker & \texttt{SIGKILL} & V3 & 2, 3 \\
F3 & Arresto del broker & \texttt{SIGKILL} & V3 & 2, 3 \\
F4 & Arresto del database (in V4: del primario) & \texttt{SIGKILL} & V1--V4 & 1, 2, 4 \\
F5 & Latenza verso il database (200~ms per 60~s) & Toxiproxy & V2, V3 & 2 \\
F6 & Perdita di un nodo \textit{(facoltativo)} & arresto del nodo & V2, V3 & 3, 6 \\
F7 & Blocco di un'istanza dell'API, con due configurazioni delle sonde & \texttt{SIGSTOP} & V2 & 3 \\
\bottomrule
\end{tabularx}
\end{table}

\paragraph{Crash e blocco.} Le due classi principali si comportano in modo opposto rispetto a Kubernetes. Un processo che \textbf{termina} viene rilevato quasi subito, perché il container esce e il kubelet lo riavvia. Un processo \textbf{bloccato} resta in esecuzione: Kubernetes se ne accorge solo dalle sonde, dopo un numero stabilito di controlli falliti di fila. Il tempo di rilevamento dipende quindi da periodo e soglia delle sonde, ed è ciò che il guasto F7 misura con due configurazioni (\emph{veloce}: controllo ogni 1~s; \emph{standard}, i valori predefiniti di Kubernetes: ogni 10~s).

La figura~\ref{fig:anatomia} scompone il fermo in rilevamento, riavvio e ritorno a pronto. La durata delle barre è indicativa.

\begin{figure}[H]
\centering
\begin{tikzpicture}[x=0.25cm, y=0.9cm, every node/.append style={font=\scriptsize},
  seg/.style={draw=white, line width=0.6pt, minimum height=0.55cm, inner sep=0pt, anchor=west},
  ril/.style={seg, fill=accento!75, text=white},
  rav/.style={seg, fill=unina, text=white},
  pro/.style={seg, fill=unina!45},
  sta/.style={seg, fill=grigiotesto!25}]
\newcommand{\segm}[5]{\node[#1, minimum width=#3*0.25cm] at (#2,#4) {#5};}
\node[anchor=east, font=\small\bfseries, text=unina, align=right] at (-0.6,0) {Crash\\[-1pt]{\scriptsize\normalfont guasto F1}};
\segm{ril}{0}{1}{0}{}
\segm{rav}{1}{6}{0}{}
\segm{pro}{7}{5}{0}{}
\segm{sta}{12}{5}{0}{}
\node[anchor=east, font=\small\bfseries, text=unina, align=right] at (-0.6,-1) {Blocco\\[-1pt]{\scriptsize\normalfont sonde veloci}};
\segm{ril}{0}{6}{-1}{}
\segm{rav}{6}{6}{-1}{}
\segm{pro}{12}{5}{-1}{}
\segm{sta}{17}{5}{-1}{}
\node[anchor=east, font=\small\bfseries, text=unina, align=right] at (-0.6,-2) {Blocco\\[-1pt]{\scriptsize\normalfont sonde standard}};
\segm{ril}{0}{30}{-2}{periodo $\times$ soglia}
\segm{rav}{30}{6}{-2}{}
\segm{pro}{36}{5}{-2}{}
\segm{sta}{41}{5}{-2}{}
\draw[vita] (0,0.55) -- (0,-2.6);
\fill[accento] (0,0.55) circle (2.2pt);
\node[text=accento, above] at (0,0.6) {guasto $t_g$};
\draw[-{Latex}, grigiotesto] (0,-2.6) -- (47,-2.6) node[right] {tempo};
\foreach \r/\e in {0/12, -1/17, -2/41} {
  \fill[accento] (\e,\r) circle (2.2pt);
  \node[text=accento, anchor=south] at (\e,\r+0.3) {$t_r$};
}
\foreach \s/\t/\x in {ril/rilevamento/0, rav/riavvio/10, pro/ritorno a pronto/20, sta/5 s di servizio stabile/32} {
  \node[\s, minimum width=0.3cm, minimum height=0.25cm] at (\x,-3.5) {};
  \node[anchor=west, etichetta] at (\x+1.5,-3.5) {\t};
}
\end{tikzpicture}
\caption{Anatomia di un fermo con una sola copia. Il ripristino $t_r$ è l'inizio dei 5 secondi di servizio stabile, che lo confermano ma non entrano nel fermo $t_r - t_g$. Con più copie (guasto F7 su V2) il fermo finisce prima, quando la sonda di readiness esclude la copia dal Service.}
\label{fig:anatomia}
\end{figure}

\paragraph{Come si provocano.} I segnali \texttt{SIGKILL} (crash) e \texttt{SIGSTOP} (blocco) sono inviati dal nodo che ospita il processo, non con \texttt{kubectl delete pod}, che è una terminazione ordinata e non un guasto (D0016). L'istante del guasto è registrato al millisecondo. Prima di ogni run il sistema viene riportato a uno stato iniziale verificato. La latenza del guasto F5 è di 200~ms: con 500~ms il sincrono fallirebbe quasi del tutto e il guasto non direbbe nulla di nuovo rispetto all'arresto del database.

\section{Come misuriamo}
\label{s:misura}

\subsection{Che cosa misuriamo}

Tempi di fermo e di ripristino si calcolano dal log del client e dagli istanti registrati nel database: la risoluzione di Prometheus (alcuni secondi) è troppo grossolana per guasti che durano pochi secondi.

\begin{table}[H]
\centering
\caption{Le misure e il loro significato.}
\label{tab:definizioni}
\begin{tabularx}{\linewidth}{@{}>{\raggedright\arraybackslash}p{2.9cm}>{\raggedright\arraybackslash}p{3.6cm}L@{}}
\toprule
\textbf{Misura} & \textbf{In parole} & \textbf{Significato} \\
\midrule
Disponibilità di accettazione & Il cliente riceve una risposta positiva? & Frazione di secondi in cui almeno il \parametro{soglia_secondo_disponibile_percentuale}\% delle richieste inviate riceve esito positivo (201 o 202) entro \parametro{timeout_client_s}~s, visto dal client. \\
Disponibilità di completamento & L'ordine è stato davvero salvato, e in tempo? & Frazione di secondi in cui almeno il \parametro{soglia_secondo_disponibile_percentuale}\% degli ordini ricevuti è salvato nel database entro \parametro{limite_completamento_s}~s dalla ricezione. In V3 e V5 è la misura decisiva: con il database fermo l'API continua ad accettare, ma nessun ordine viene completato. \\
Ripristino $t_r$ & Da quando il servizio funziona di nuovo, in modo stabile & Vale separatamente per accettazione e completamento: primo istante, dal guasto in poi, da cui seguono 5 secondi consecutivi disponibili in quella misura. Un pod tornato pronto non basta. Il servizio è ripristinato quando lo sono sia l'accettazione sia il completamento. \\
Fermo $D$ & Quanto è durato il disservizio & $t_r - t_g$, con $t_g$ istante del guasto, calcolato con il $t_r$ di ciascuna misura. Se il ripristino non avviene entro il tetto del run, il run è valido e registrato come ``non ripristinato''. \\
Tempo end-to-end & Quanto passa dall'arrivo dell'ordine al suo salvataggio & Scrittura nel database meno ricezione da parte dell'API, entrambi presi nel cluster. \\
Costo & Quante risorse abbiamo usato & In locale, \emph{worker-secondi}: la somma nel tempo del numero di worker accesi (RQ5). Su AWS, costo per mille ordini. \\
\bottomrule
\end{tabularx}
\end{table}

La soglia del \parametro{soglia_secondo_disponibile_percentuale}\% è severa a ritmi bassi: a 50 richieste al secondo basta un errore per rendere indisponibile l'intero secondo. Per questo si riportano sempre anche le \emph{richieste fallite per guasto}, che non dipendono dalla soglia, e nell'analisi il limite di completamento si ripete con 1~s e 5~s come verifica di sensibilità.

\subsection{Dove finisce ogni ordine}

A fine run ogni identificativo inviato dal client viene assegnato a una categoria, incrociando il log del client con il database e la coda dei falliti (figura~\ref{fig:categorie}). Il caso grave è il \textbf{perso}: il client ha ricevuto esito positivo, ma l'ordine non è mai stato salvato. Distinguere \emph{fallito} (errore dichiarato, il client riprova) da \emph{perso} è ciò che permette di confrontare sincrono e asincrono: nel sincrono un errore è visibile subito, nell'asincrono un ``ricevuto'' non garantisce il completamento. I \textbf{doppi} (lo stesso ordine consegnato due volte) si contano a parte: l'identificativo dell'ordine lo sceglie il client e il database inserisce un ordine solo se non esiste già, registrando i doppi in una tabella separata.

\begin{figure}[H]
\centering
\resizebox{\linewidth}{!}{%
\begin{tikzpicture}[
  domanda/.style={draw=unina, thick, fill=white, rounded corners=8pt, align=center, font=\scriptsize, minimum height=0.85cm, inner sep=3pt},
  cat/.style={box, minimum width=3.3cm, minimum height=0.8cm, font=\scriptsize},
  si/.style={etichetta, above, pos=0.45},
  no/.style={etichetta, below, pos=0.45}]
\node[domanda] (q1) at (0,0) {Il client ha ricevuto\\esito positivo?};
\node[domanda] (q2) at (4.4,1.5) {L'ordine è\\nel database?};
\node[domanda] (q3) at (4.4,-1.5) {L'ordine è\\nel database?};
\node[domanda] (q4) at (8.4,0.6) {È nella coda\\dei falliti?};
\node[cat] (c1) at (8.4,2.35) {\textbf{completato in tempo}};
\node[cat, fill=uninachiaro!40] (c2) at (12.3,1.25) {\textbf{in dead-letter}\\rielaborato e contato a parte};
\node[cat, guastato] (c3) at (12.3,-0.05) {\textbf{perso}\\il caso grave};
\node[cat, fill=white] (c4) at (8.4,-1.0) {\textbf{completato dopo il timeout}};
\node[cat, fill=white] (c5) at (8.4,-2.3) {\textbf{fallito}\\errore dichiarato};
\draw[freccia] (q1) |- node[si, pos=0.75] {sì} (q2);
\draw[freccia] (q1) |- node[no, pos=0.75] {no (errore o timeout)} (q3);
\draw[freccia] (q2) |- node[si, pos=0.7] {sì} (c1);
\draw[freccia] (q2) |- node[no, pos=0.7] {no} (q4);
\draw[freccia] (q4) |- node[si, pos=0.7] {sì} (c2);
\draw[freccia] (q4) |- node[no, pos=0.7] {no} (c3);
\draw[freccia] (q3) |- node[si, pos=0.7] {sì} (c4);
\draw[freccia] (q3) |- node[no, pos=0.7] {no} (c5);
\end{tikzpicture}}
\caption{Assegnazione di un ordine a una categoria. I duplicati sono contati a parte.}
\label{fig:categorie}
\end{figure}

\subsection{Un run}

Ogni run è automatico (figura~\ref{fig:run}): reset verificato (coda vuota, tabelle ripulite, pod nello stato iniziale), un minuto di riscaldamento escluso dalle statistiche, un minuto di riferimento senza guasti, guasto al minuto 2, osservazione fino al ripristino più 2 minuti, ed esportazione dei dati. Il tetto è di \parametro{tetto_min} minuti, \parametro{tetto_min_f6_e_w2} per il guasto F6 (Kubernetes attende alcuni minuti prima di ricollocare i pod di un nodo perso) e per i run sotto W2 (con un solo worker la coda si svuota in circa 8 minuti). I run senza guasti durano 5 minuti fissi.

Il carico segue un \textbf{modello aperto}: le richieste partono a ritmo fisso anche quando il sistema rallenta. Con un modello chiuso il generatore smetterebbe di inviare proprio durante il guasto, nascondendo il problema (\textit{coordinated omission}).

\begin{figure}[H]
\centering
\begin{tikzpicture}[x=1.2cm, every node/.append style={font=\scriptsize}]
\fill[grigiotesto!20] (-1,0) rectangle (0,0.7);
\node at (-0.5,0.35) {reset};
\fill[unina!25] (0,0) rectangle (1,0.7);
\node[font=\tiny, align=center] at (0.5,0.35) {riscal-\\damento};
\fill[unina!55] (1,0) rectangle (2,0.7);
\node[text=white, font=\tiny, align=center] at (1.5,0.35) {riferi-\\mento};
\fill[unina!85] (2,0) rectangle (6,0.7);
\node[text=white, align=center] at (4,0.35) {osservazione\\fino al ripristino + 2 min};
\fill[unina!10] (6,0) rectangle (10,0.7);
\draw[unina!60, densely dashed] (6,0) rectangle (10,0.7);
\node at (8,0.35) {solo se il ripristino tarda};
\fill[grigiotesto!20] (10,0) rectangle (11.2,0.7);
\node at (10.6,0.35) {export};
\draw[-{Latex}, grigiotesto] (-1,-0.15) -- (11.4,-0.15);
\foreach \t/\l in {0/0, 1/1, 2/2, 4/4, 6/6, 8/8, 10/10} {
  \draw[grigiotesto] (\t,-0.1) -- (\t,-0.2);
  \node[below, text=grigiotesto] at (\t,-0.2) {\l};
}
\node[below, text=grigiotesto, anchor=north west] at (11.2,-0.2) {min};
\draw[accento, very thick] (2,-0.25) -- (2,1.05);
\node[text=accento, above, align=center] at (2,1.05) {guasto $t_g$\\(minuto 2)};
\draw[accento, densely dashed, thick] (10,-0.25) -- (10,1.05);
\node[text=accento, above, align=center] at (10,1.05) {tetto: \parametro{tetto_min} min\\(\parametro{tetto_min_f6_e_w2} per F6 e W2)};
\draw[decorate, decoration={brace, amplitude=4pt, mirror}, grigiotesto] (0,-0.75) -- node[below=3pt, etichetta] {escluso dalle statistiche} (1,-0.75);
\end{tikzpicture}
\caption{Struttura di un run: la durata tipica è di 5--6 minuti.}
\label{fig:run}
\end{figure}

\subsection{Carichi e numero di run}

\begin{table}[H]
\centering
\caption{I tre profili di carico. I livelli si fissano rispetto alla capacità misurata nella fase di banco di prova (F4).}
\label{tab:carichi}
\begin{tabularx}{\linewidth}{@{}llL@{}}
\toprule
\textbf{ID} & \textbf{Profilo} & \textbf{Uso} \\
\midrule
W1 & costante, sotto la capacità & Tutti i run con guasto e i run di riferimento. Ritmo uguale per tutte le versioni: il \parametro{w1_percentuale_capacita_minima}\% della capacità della versione più lenta, così il carico non diventa un secondo fattore \\
W2 & picco improvviso, poi ritorno alla base & Elasticità (RQ5). Base al \parametro{w2_base_percentuale_mu1}\% della capacità $\mu_1$ di un worker, picco al \parametro{w2_picco_percentuale_mu1}\% di $\mu_1$ per \parametro{w2_durata_picco_s}~s dal minuto 2: sopra 3 worker e sotto 10 \\
W3 & sovraccarico prolungato & Comportamento oltre il limite, senza guasti, su V1, V2 e V3: \parametro{w3_percentuale_capacita}\% della capacità di ciascuna \\
\bottomrule
\end{tabularx}
\end{table}

\begin{table}[H]
\centering
\caption{Numero di run previsti.}
\label{tab:matrice}
\begin{tabular}{@{}lrr@{}}
\toprule
\textbf{Gruppo} & \textbf{Ripetizioni} & \textbf{Run} \\
\midrule
Senza guasti, carico W1 (V1, V2, V3, V4a, V4s) & 5 & 25 \\
Senza guasti, carico W3 (V1, V2, V3) & 5 & 15 \\
Con guasti F1--F5 su V1--V3 & 10 & 100 \\
Guasto del database su V4 (V4a, V4s) & 10 & 20 \\
Blocco dell'API (F7, due configurazioni delle sonde) & 10 & 20 \\
Carico W2: V5 e V3 con 3, con 1 e con 10 worker fissi & 10 & 40 \\
\midrule
\textbf{Totale} & & \textbf{\parametro{run_totali}} \\
\bottomrule
\end{tabular}
\end{table}

Il confronto di RQ5 riguarda quattro configurazioni sotto W2, tutte con 10 ripetizioni per avere la stessa precisione. I run senza guasti danno il quadro di base di ogni versione e confermano che, senza guasti, il fermo è zero. La figura~\ref{fig:matrice} mostra come i 180 run con guasto o sotto W2 si distribuiscono; il guasto F6 è facoltativo e non è nel totale.

\begin{figure}[H]
\centering
\begin{tikzpicture}[x=1.55cm, y=0.62cm, every node/.append style={font=\scriptsize}]
\foreach \c/\v in {1/V1, 2/V2, 3/V3, 4/V4a, 5/V4s, 6/V5} \node[font=\scriptsize\bfseries, text=unina] at (\c,0.9) {\v};
\foreach \r/\f in {1/F1 crash API, 2/F2 crash worker, 3/F3 crash broker, 4/F4 crash database, 5/F5 latenza DB, 6/F6 perdita nodo, 7/F7 blocco API, 8/V5 e riferimenti (W2)}
  \node[anchor=east] at (0.45,-\r+1) {\f};
\foreach \c in {1,...,6} \foreach \r in {1,...,8}
  \draw[white, fill=black!4] (\c-0.47,-\r+1-0.45) rectangle (\c+0.47,-\r+1+0.45);
\newcommand{\cella}[4]{\draw[white, fill=#3] (#1-0.47,-#2+1-0.45) rectangle (#1+0.47,-#2+1+0.45); \node at (#1,-#2+1) {#4};}
\cella{1}{1}{unina!30}{10}\cella{2}{1}{unina!30}{10}\cella{3}{1}{unina!30}{10}
\cella{3}{2}{unina!30}{10}
\cella{3}{3}{unina!30}{10}
\cella{1}{4}{unina!30}{10}\cella{2}{4}{unina!30}{10}\cella{3}{4}{unina!30}{10}\cella{4}{4}{unina!30}{10}\cella{5}{4}{unina!30}{10}
\cella{2}{5}{unina!30}{10}\cella{3}{5}{unina!30}{10}
\cella{2}{6}{black!12}{opz.}\cella{3}{6}{black!12}{opz.}
\cella{2}{7}{unina!55}{\textcolor{white}{2 $\times$ 10}}
\cella{3}{8}{unina!55}{\textcolor{white}{3 $\times$ 10}}\cella{6}{8}{unina!55}{\textcolor{white}{10}}
\end{tikzpicture}
\caption{Run per guasto e versione (140 con guasto più 40 sotto W2). Ogni cella vale 10 ripetizioni; per F7 le due configurazioni delle sonde, sotto W2 V3 con 3, con 1 e con 10 worker fissi. I 40 run senza guasti non sono nella figura.}
\label{fig:matrice}
\end{figure}

I run si eseguono in \textbf{ordine casuale}, mescolando configurazioni e ripetizioni, così un eventuale degrado dell'ambiente non si concentra su una sola configurazione. La campagna richiede circa \parametro{ore_campagna_circa} ore di macchina.

\subsection{Statistica e validità}

Per ogni configurazione si riportano \textbf{mediana} e \textbf{intervallo interquartile}; con 10 ripetizioni anche l'\textbf{intervallo di confidenza al 95\%} della mediana, calcolato con il bootstrap. Con 5 ripetizioni non si traggono conclusioni su differenze piccole, e per ogni confronto si dichiara la differenza minima distinguibile.

Un run è valido se il reset è verificato, se il ritmo inviato resta entro $\pm 5\%$ di quello previsto e se la CPU e la memoria del computer restano entro le soglie fissate nel banco di prova. Un run non valido viene scartato e ripetuto, annotando il motivo nel diario.

\section{Che cosa ci aspettiamo}
\label{s:previsioni}

Le previsioni si scrivono e si registrano su Git \textbf{prima} di qualsiasi misura, con il tag \texttt{m1}: la data del commit dimostra che non sono state adattate ai risultati. Dopo il tag non si modificano; uno scarto tra previsto e misurato si spiega nel report, non si corregge.

\subsection{Le grandezze}

\begin{table}[H]
\centering
\caption{Le grandezze usate nelle previsioni, con un esempio.}
\label{tab:grandezze}
\begin{tabularx}{\linewidth}{@{}lLl@{}}
\toprule
\textbf{Simbolo} & \textbf{Significato} & \textbf{Esempio} \\
\midrule
$\lambda$ & ritmo di arrivo: ordini inviati al secondo & 100 ordini/s \\
$\mu$ & capacità: ordini al secondo che il sistema smaltisce & 200 ordini/s \\
$\rho = \lambda / \mu$ & utilizzazione: quanto il sistema è occupato & 50\% \\
$N$ & copie dell'API & 3 \\
$D = t_r - t_g$ & fermo: per quanto il servizio resta indisponibile & 4 s \\
\bottomrule
\end{tabularx}
\end{table}

\paragraph{Una sola copia, sincrono (guasto F1 su V1).} Durante il fermo fallisce tutto: $\text{falliti} \approx \lambda \cdot D = 100 \cdot 4 = 400$ ordini.

\paragraph{Tre copie (guasto F1 su V2).} Fallisce solo il traffico diretto alla copia guasta, finché non viene esclusa dal Service ($T_{\text{escl}} \approx 1{,}5$~s): $\text{falliti} \approx \frac{\lambda}{N} T_{\text{escl}} = \frac{100}{3} \cdot 1{,}5 \approx 50$ ordini. Il fermo non è nullo: per quei secondi fallisce un terzo delle richieste, molto più dell'1\% tollerato.

\paragraph{Con la coda (guasto F4 su V3).} Il database resta fermo $D_0 = 30$~s: l'API continua ad accettare, ma gli ordini si accumulano ($\lambda D_0 = 3000$). Quando il database torna, la coda si svuota alla velocità netta $\mu - \lambda = 100$ ordini/s:
\[
T_{\text{arretrato}} \approx \frac{\lambda \cdot D_0}{\mu - \lambda} = \frac{3000}{100} = 30 \text{ s.}
\]
Nessun ordine fallisce, ma il completamento torna regolare solo dopo $D_0 + T_{\text{arretrato}} = 60$~s. Per questo nell'asincrono si prevedono separatamente il fermo di accettazione e quello di completamento.

\subsection{Il modello del fermo}

Il fermo previsto si ricava dai parametri configurati, non a occhio:
\begin{equation}
  D \approx \underbrace{T_{\text{ril}}}_{\text{rilevamento}} + \underbrace{T_{\text{riavvio}}}_{\text{avvio del container}} + \underbrace{T_{\text{pronto}}}_{\text{prima readiness positiva}} + \underbrace{T_{\text{arretrato}}}_{\text{solo V3, V5}}
  \label{eq:fermo}
\end{equation}
Per un crash il rilevamento è quasi nullo; per un blocco vale circa $p \times s + t$, con $p$ periodo della sonda, $s$ soglia di fallimenti consecutivi e $t$ timeout. In V4 il riavvio è sostituito dalla promozione dello standby.

Alcuni tempi si conosceranno solo con la misura a vuoto e i primi test del banco di prova. Per scrivere le previsioni prima, si fissano valori \emph{ipotizzati}, uguali per tutte le righe (tabella~\ref{tab:ipotesi}). Se un valore misurato esce dall'intervallo, la previsione non si corregge: lo scarto si spiega nel report.

\begin{table}[H]
\centering
\caption{I principali valori ipotizzati. L'elenco completo è nel documento tecnico.}
\label{tab:ipotesi}
\begin{tabularx}{\linewidth}{@{}lLl@{}}
\toprule
\textbf{Simbolo} & \textbf{Cosa} & \textbf{Valore} \\
\midrule
$T_{\text{ril}}$ & rilevamento di un crash da parte del kubelet & $\approx 1$~s \\
$T_{\text{riavvio}}$ & riavvio del container dell'API o del worker & 1--3~s \\
$T_{\text{escl}}$ & uscita dal Service di una copia dell'API in crash & 1--2~s \\
$T_{\text{DB}}$ & riavvio di PostgreSQL dopo \texttt{SIGKILL} & 3--8~s \\
$T_{\text{MQ}}$ & riavvio di RabbitMQ con la coda durevole & 5--15~s \\
$T_{\text{prom}}$ & promozione dello standby da parte di CloudNativePG & 3--10~s \\
$T_{\text{KEDA}}$ & reazione di V5 al picco (controllo ogni 15~s, nuovi worker pronti in 2--5~s) & 3--20~s i primi worker; 18--35~s tutti e 10 \\
$\mu_N$ & capacità con $N$ worker; se è più bassa, nel report si indica quale componente si è saturato & $N\mu_1$, fino a 10 worker \\
\bottomrule
\end{tabularx}
\end{table}

\subsection{Le previsioni}

Nell'asincrono il fermo si prevede separatamente per l'accettazione (acc.) e per il completamento (compl.). I numeri tra parentesi valgono per $\lambda = 100$ ordini/s e sono solo un esempio: il valore vero di $\lambda$ si fissa nel banco di prova.

\begin{table}[H]
\centering
\caption{Previsioni per guasto e versione.}
\label{tab:previsioni}
\small
\begin{tabularx}{\linewidth}{@{}llLLcc@{}}
\toprule
\textbf{Guasto} & \textbf{Vers.} & \textbf{Fermo previsto} & \textbf{Falliti} & \textbf{Persi} & \textbf{Doppi} \\
\midrule
F1 & V1 & 3--5~s & $\lambda D$ ($\approx$ 300--500) & 0 & 0 \\
F1 & V2 & 1--2~s & $\frac{\lambda}{3} T_{\text{escl}}$ ($\approx$ 30--70) & 0 & 0 \\
F1 & V3 & acc.\ 1--2~s; compl.\ $\approx 0$ & $\frac{\lambda}{3} T_{\text{escl}}$ ($\approx$ 30--70) & 0 & 0 \\
F2 & V3 & acc.\ 0; compl.\ 6--18~s & 0 & 0 & 0--3 \\
F3 & V3 & acc.\ e compl.\ 7--21~s & $\lambda D$ ($\approx$ 700--2100) & 0 & 0--3 \\
F4 & V1, V2 & 5--10~s & $\lambda D$ ($\approx$ 500--1000) & 0 & 0 \\
F4 & V3 & acc.\ 0; compl.\ 10--36~s & 0 & 0 & 0 \\
F4 & V4a & 5--12~s & $\lambda D$ ($\approx$ 500--1200) & $\le 0{,}1\lambda$ ($\approx$ 10) & 0 \\
F4 & V4s & 5--12~s & $\lambda D$ ($\approx$ 500--1200) & 0 & 0 \\
F5 & V2 & 60~s (tutta la durata) & $(\lambda - \mu_L) \cdot 60$ & 0 & 0 \\
F5 & V3 & acc.\ 0; compl.\ $\approx$ 120~s & 0 & 0 & 0 \\
F7 & V2 veloce & $\approx$ 4~s & $\frac{\lambda}{3} D$ ($\approx$ 130) & 0 & 0 \\
F7 & V2 standard & $\approx$ 31~s & $\frac{\lambda}{3} D$ ($\approx$ 1000) & 0 & 0 \\
\bottomrule
\end{tabularx}
\end{table}

Due derivazioni mostrano come si ottengono i valori. Nel guasto F7 su V2 un processo bloccato non risponde nemmeno con un errore: il fermo dura finché la sonda di readiness non esclude la copia, cioè $p \times s + t = 1 \times 3 + 1 \approx 4$~s con le sonde veloci e $10 \times 3 + 1 \approx 31$~s con quelle standard. Nel guasto F5 la latenza $L = 200$~ms si somma a ognuno dei $k = 2$--$3$ scambi con il database, quindi ogni ordine richiede circa 0,4--0,6~s e il database regge al massimo $\mu_L = C / (S + kL)$ ordini al secondo, con $C$ scritture contemporanee: in V2 la parte in eccesso fallisce per tutti i 60~s, in V3 si accumula in coda e si smaltisce poi in un tempo simile alla durata del guasto.

\subsection{Le relazioni attese}

I numeri dipendono dalle ipotesi; le relazioni tra versioni no, e sono la parte più facile da smentire.

\begin{table}[H]
\centering
\caption{Relazioni attese per domanda di ricerca.}
\label{tab:relazioni}
\begin{tabularx}{\linewidth}{@{}lL@{}}
\toprule
\textbf{RQ} & \textbf{Che cosa prevediamo} \\
\midrule
RQ1 & Con F1, V2 ha un fermo più breve di V1 e circa un ordine di grandezza in meno di ordini falliti. Con F4 i due si comportano allo stesso modo: il database resta unico. \\
RQ2 & Con F4, V3 non perde accettazione e ha zero falliti, ma il suo fermo di completamento è più lungo di quello di V2. Nessun ordine si perde, tranne in V4a. Il broker è il nuovo punto debole: F3 su V3 costa più di F4 su V2. \\
RQ3 & Con F7 le sonde standard allungano il fermo di circa 8 volte rispetto alle veloci. Un crash (F1) si rileva in circa 1~s, un blocco (F7) solo con le sonde. \\
RQ4 & In locale il failover di V4 non è necessariamente più rapido del riavvio del database di V2, perché il riavvio dura pochi secondi: il vantaggio sarebbe nella perdita di un nodo, che qui non si prova. V4a può perdere pochi ordini, V4s nessuno, al prezzo di scritture più lente. \\
RQ5 & Sotto W2, V5 sta tra V3 a 1 e V3 a 10 worker sia per arretrato massimo sia per costo. V5 impiega circa 18--35~s per arrivare a 10 worker; nel frattempo la coda cresce di $4\mu_1$ ordini al secondo, fino a circa $80\mu_1$. V3 con i suoi 3 worker ha capacità $3\mu_1$ contro $5\mu_1$ in arrivo: dopo 60~s la coda arriva a circa $120\mu_1$, più di V5. \\
RQ6 & Su AWS il fermo di F1 su V2 è più lungo che in locale, perché l'ALB toglie una copia con controlli di salute più lenti di quelli del Service. I risultati che non dipendono dai tempi (zero ordini persi in V3) si ripetono uguali. \\
\bottomrule
\end{tabularx}
\end{table}

\subsection{Dal fermo alla disponibilità}

Un fermo di 4~s, da solo, non dice se il servizio è buono: dipende da quante volte il guasto si ripete. Portiamo i secondi misurati sulla scala dei servizi reali, la disponibilità del 99,9\%. Un mese di 30 giorni ha 43\,200 minuti, quindi il \textbf{budget di errore} è
\[
B = 0{,}001 \cdot 43\,200 \text{ min} = 43{,}2 \text{ min} = 2592 \text{ s al mese.}
\]
Quante volte al mese i guasti avverranno davvero non si può misurare, perché li provochiamo noi. Si rovescia allora la domanda: quanti guasti di un certo tipo può sopportare una versione prima di esaurire il budget? Sono $B / D$ (tabella~\ref{tab:sopportabili}); nel report il calcolo si ripete con i fermi misurati.

\begin{table}[H]
\centering
\caption{Guasti sopportabili al mese con l'obiettivo del 99,9\%, calcolati sui fermi previsti.}
\label{tab:sopportabili}
\begin{tabular}{@{}lrr@{}}
\toprule
\textbf{Guasto} & \textbf{Fermo $D$} & \textbf{Guasti sopportabili al mese} \\
\midrule
F1 su V1 & 4~s & $\approx 650$ \\
F1 su V2 & 1,5~s & $\approx 1700$ \\
F4 su V1 e V2 & 8~s & $\approx 320$ \\
F7, sonde veloci & 4~s & $\approx 650$ \\
F7, sonde standard & 31~s & $\approx 80$ \\
\bottomrule
\end{tabular}
\end{table}

Con le sonde standard bastano circa 80 blocchi al mese (meno di 3 al giorno) per esaurire il budget; con quelle veloci ne servono circa 650.

\section{Le decisioni principali}
\label{s:decisioni}

Ogni scelta rilevante è documentata in una scheda con contesto, decisione e conseguenze (ADR) nel documento tecnico. La tabella~\ref{tab:decisioni} ne riassume l'essenza: che cosa abbiamo scelto e perché.

\begin{table}[H]
\centering
\caption{Le decisioni architetturali in breve.}
\label{tab:decisioni}
\small
\begin{tabularx}{\linewidth}{@{}>{\raggedright\arraybackslash}p{1.9cm}LL@{}}
\toprule
\textbf{Decisione} & \textbf{Che cosa abbiamo scelto} & \textbf{Perché} \\
\midrule
D0001, D0002 & Con la coda l'API risponde 202 solo dopo che RabbitMQ ha confermato di aver scritto il messaggio (durevole); il worker conferma dopo il salvataggio; un ordine alla volta per worker; coda dei falliti & Un ordine ``ricevuto'' non deve sparire se coda o worker si spengono. Prezzo: l'API è più lenta e un ordine può arrivare due volte \\
D0003 & Il worker ritenta tenendo l'ordine con sé, con attese crescenti da 0,5 a 10~s; dopo 120~s va alla coda dei falliti & Rimetterlo subito in coda creerebbe un giro continuo senza risultato. I 120~s sono indicativi \\
D0004 & Identificativo scelto dal client; inserimento ``se esiste non fare nulla''; doppi registrati in una tabella & I doppi si contano anche quando il programma si spegne \\
D0005 & Ricezione e scrittura prese entrambe dal cluster, mai dal client & Orologi diversi falserebbero il tempo di ogni ordine \\
D0006 & Database e coda su un nodo dedicato con regola rigida; primario e standby su nodi diversi; API e worker preferiscono l'altro nodo & Collocazione ripetibile: se cambiasse tra i run sarebbe una causa nascosta delle differenze \\
D0007 & Le sonde controllano solo l'API, non il database & Un database fermo non deve far riavviare o escludere tutte le copie, annullando il vantaggio dell'asincrono \\
D0008 & Il generatore di carico aspetta 2~s, l'API 1,5~s e poi risponde 503; limite di completamento di 2~s & Un errore dichiarato è più chiaro di un timeout, e l'API non lavora per chi ha già rinunciato \\
D0010 & Pool di 5 connessioni per pod, uguale in tutte le versioni (V5 al massimo 50 su 100) & Il limite di PostgreSQL non deve falsare il confronto \\
D0014 & V4 con CloudNativePG, in due modalità: asincrona (V4a) e sincrona (V4s) & In V1--V3 il database è unico: se cade si ferma tutto \\
D0015 & V5 con KEDA da 1 a 10 worker: controllo ogni 15~s circa, soglia di due secondi di lavoro in coda per worker, attesa di 60~s prima di toglierli & Il valore predefinito (300~s) terrebbe 10 worker accesi dopo il picco e annullerebbe il confronto di costo \\
D0016 & Crash e blocco con \texttt{SIGKILL} e \texttt{SIGSTOP} inviati dal nodo & \texttt{kubectl delete} è una chiusura ordinata, non un guasto \\
D0017 & Obiettivo di disponibilità del 99,9\% (circa 43 minuti di fermo al mese) & Scala comune per confrontare versioni e guasti \\
D0009, D0011--D0013 & AWS in Irlanda con \texttt{eksctl} e istanze a prestazioni fisse; Python; un solo codice per tutte le versioni; misure ufficiali su un solo computer da 16 GB & Confronto equo e risorse stabili tra i run \\
\bottomrule
\end{tabularx}
\end{table}

\section{I limiti dell'esperimento}
\label{s:limiti}

Nessun esperimento è privo di limiti. Le minacce principali sono classificate nelle quattro categorie usuali, ognuna con la sua contromisura. \emph{Interna}: un altro fattore causa la differenza osservata. \emph{Di costrutto}: la misura non rappresenta ciò che vogliamo misurare. \emph{Esterna}: i risultati non valgono fuori dall'ambiente di prova. \emph{Delle conclusioni}: la statistica porta a conclusioni sbagliate.

\begin{table}[H]
\centering
\caption{Le minacce principali e le contromisure.}
\label{tab:limiti}
\small
\begin{tabularx}{\linewidth}{@{}>{\raggedright\arraybackslash}p{1.8cm}LL@{}}
\toprule
\textbf{Tipo} & \textbf{Minaccia} & \textbf{Contromisura} \\
\midrule
Interna & Sistema e generatore di carico girano sullo stesso computer e si contendono CPU e memoria & Soglia di CPU fissata nel banco di prova; run saturi scartati \\
Interna & Un run lascia tracce (code non vuote, pod degradati) che influenzano il successivo & Reset verificato prima di ogni run; ordine casuale dei run \\
Di costrutto & Un generatore a modello chiuso smetterebbe di inviare durante il guasto & Carico a modello aperto; ritmo controllato entro $\pm 5\%$ \\
Di costrutto & Un pod ``pronto'' non significa servizio ripristinato & Ripristino come 5~s consecutivi disponibili in accettazione e completamento \\
Esterna & In locale i nodi sono container sullo stesso hardware: la perdita di un nodo è simulata & Parte degli esperimenti ripetuta su AWS (RQ6); la perdita di un nodo solo se si esegue il guasto F6, facoltativo \\
Esterna & Il servizio è molto più semplice di uno reale & Le conclusioni riguardano i meccanismi (rilevamento, riparazione, smaltimento), non i valori assoluti \\
Esterna & Il control plane di Kubernetes è unico e funge anche da punto d'ingresso & Non viene mai guastato; in produzione sarebbe replicato \\
Delle conclusioni & Un risultato singolo può essere casuale & 10 ripetizioni per i run con guasto; intervalli di confidenza bootstrap \\
Delle conclusioni & Si possono adattare le spiegazioni dopo aver visto i dati & Previsioni registrate prima delle misure (tag \texttt{m1}) \\
\bottomrule
\end{tabularx}
\end{table}

\paragraph{Fuori ambito.} Replica del broker e del control plane, database distribuito su più zone, sicurezza e ottimizzazione delle prestazioni del codice.

\section{Piano di lavoro}
\label{s:piano}

Il progetto dura dodici settimane, dal 22/10/2026 al 13/01/2027, per circa 306 ore in due. Ogni milestone è un controllo: se non è superato, non si va avanti.

\begin{table}[H]
\centering
\caption{Le milestone.}
\label{tab:milestone}
\begin{tabularx}{\linewidth}{@{}llL@{}}
\toprule
\textbf{Milestone} & \textbf{Entro il} & \textbf{Criterio di passaggio} \\
\midrule
M1 & 04/11/2026 & Decisioni scritte, parametri fissati, previsioni registrate (tag \texttt{m1}) \\
M2 & 11/11/2026 & Applicazione minima V1--V3 funzionante, nessun ordine perso con il database fermo per 60~s \\
M3 & 25/11/2026 & V1--V3 su kind a 3 nodi; un crash (guasto F1) misurato con risoluzione sotto il secondo \\
M3b & 02/12/2026 & Account AWS con allarme di budget attivo; cluster creato e distrutto, costi e limiti noti \\
M4 & 09/12/2026 & Banco di prova ripetibile; failover di V4 e scalatura di V5 verificati \\
M5 & 16/12/2026 & 220 run validi; decisione sull'estensione della fase AWS \\
M6 & 30/12/2026 & Sessioni AWS svolte, risorse distrutte, costi annotati \\
M7 & 13/01/2027 & Report, diario, repository e orale pronti \\
\bottomrule
\end{tabularx}
\end{table}

\paragraph{Piano di taglio.} In caso di ritardo si rinuncia, in quest'ordine, a: perdita di un nodo (guasto F6), V5, latenza del database (guasto F5), parte delle sessioni AWS, replica sincrona di V4 (V4s). Non si rinuncia mai al protocollo di misura, a V1--V3, ai guasti F1--F4 e F7 con 10 ripetizioni, alle previsioni e al diario.

\paragraph{Rischi.} Le settimane S9--S11 cadono durante le festività: sessioni AWS e analisi si confermano o si ripianificano a M5. Il rischio principale su AWS è lasciare risorse accese, per questo ogni sessione termina con la distruzione del cluster e la verifica che non restino risorse a pagamento.

\paragraph{Gestione.} Ogni attività parte da una issue assegnata a una milestone; le modifiche rientrano in \texttt{main} con una pull request rivista dall'altro autore; ogni operazione è annotata nel diario (\texttt{report/diario/}).


\end{document}
